\section{Method}
\label{sec:method}

\subsection{A multi-factor memory value}
\label{sec:method:value}

Each memory $m$ (a stored turn or consolidated item) is summarised by a
factor vector $\fvec(m)\in[0,1]^7$ and scored by a learned linear value
\begin{equation}
V(m) \;=\; \sum_{i=1}^{7} w_i\, f_i(m),
\qquad \wvec \in \mathbb{R}_{\ge 0}^{7}.
\label{eq:value}
\end{equation}
The seven factors are chosen so that each is an interpretable,
independently computable determinant of expected future usefulness
(\cref{tab:factors}). We regard \cref{eq:value} as a \emph{learned reward
model for memory}: $V(m)$ estimates a memory's contribution to future
task return, and its weights are fit by policy optimisation
(\cref{sec:method:learn}) rather than set by hand.

Linearity here is a deliberate design choice, not a limitation we
tolerate, for three reasons. \textbf{(i) Capacity match}: with seven
interpretable factors and limited, indirect supervision, a
seven-parameter value is the appropriate capacity; richer function
classes overfit this regime (\cref{sec:limitations}). \textbf{(ii)
Auditability}: the learned weights $w_i$ \emph{are} the explanation --- a
deployment can read off which factors a given workload rewards
(\cref{sec:exp:blind}), which a black-box scorer cannot provide and which
is a practical requirement for a memory controller. \textbf{(iii)
Decision-theoretic reading}: \cref{eq:value} is the first-order
(additive-utility) approximation of a memory's expected contribution to
return --- the principled starting point before modelling factor
interactions. More generally the value is any learned scoring function
$V(m)=g_\theta(\fvec(m))$, and \cref{eq:value} is its interpretable linear
instance; a neural $g_\theta$ is validated as an \emph{interaction
ablation} in \cref{sec:exp:blind} --- an MLP over the same factors ties
the linear model, confirming they combine near-additively, so the linear
value stays the default.

\begin{table}[ht!]
\centering\small
\caption{The seven memory-value factors. The right column marks which
factors are populated by the \emph{API-free} annotator used in our
experiments; the remaining three require a value profile, an LLM judge,
or access logs and are held at $0$ here (\cref{sec:limitations}).}
\label{tab:factors}
\begin{tabular}{llc}
\toprule
Factor & Definition (API-free form) & Live here? \\
\midrule
emotional intensity & $|\valence|\cdot\arousal$ from a lexical/structural extractor & yes \\
goal relevance      & $\tfrac12(1+\cos(\mathbf{e}_m, \mathbf{e}_{\text{goal}}))$ & yes \\
self/user relevance & $\tfrac12(1+\cos(\mathbf{e}_m, \mu_{\text{user}}))$ & yes \\
reliability         & provenance heuristic (user-stated $>$ model-stated) & yes \\
value alignment     & match to a value profile (SOUL) & no ($0$) \\
task utility        & marginal task-success contribution (LLM judge) & no ($0$) \\
usage history       & $r/(1+r)$, $r=$ retrieval count & no ($0$) \\
\bottomrule
\end{tabular}
\end{table}

Here $\mathbf{e}_m$ is a sentence embedding of the memory text
(all-MiniLM-L6-v2, $384$-dim; \citealp{reimers2019sbert}),
$\mu_{\text{user}}$ is the running centroid of the user's turns (a
query-agnostic ``what this user is about'' anchor), and
$\mathbf{e}_{\text{goal}}$ is the embedding of the agent's active goal.
The goal anchor is the pivot of our central experiment
(\cref{sec:exp:blind}): at retrieval time it is the query, but at
\emph{consolidation} time the agent must use a query-agnostic proxy.

\subsection{One value, three operations}
\label{sec:method:control}

The single scalar $V(m)$ drives all three memory operations
(\cref{fig:schematic}), replacing the separate, hand-tuned rules of prior
systems.

\begin{figure}[ht!]
\centering
\begin{tikzpicture}[
  font=\small, >=Stealth, node distance=0.45cm and 1.15cm,
  box/.style={draw, rounded corners, align=left, inner sep=4pt},
  val/.style={draw, rounded corners, fill=orange!12, align=center,
              inner sep=5pt, minimum height=1.0cm},
  op/.style={draw, rounded corners, fill=green!8, align=center,
             inner sep=4pt, minimum width=3.0cm}]
\node[box, fill=blue!6] (f)
  {emotion\\ goal rel.\\ value align.\\ self/user rel.\\
   task util.\\ reliability\\ usage};
\node[above=0.05cm of f, font=\footnotesize] {factors $\fvec(m)$};
\node[val, right=of f] (v) {$V(m)$\\[2pt]$=\sum_i w_i f_i$};
\node[op, right=1.2cm of v] (forget) {forget score (Eq.~\ref{eq:forget})};
\node[op, above=of forget] (enc) {encoding depth (4 tiers)};
\node[op, below=of forget] (ret) {retrieval rank};
\draw[->, thick] (f) -- (v);
\draw[->] (v.east) -- (enc.west);
\draw[->] (v.east) -- (forget.west);
\draw[->] (v.east) -- (ret.west);
\end{tikzpicture}
\caption{\textbf{One value, three operations.} A single learned scalar
$V(m)=\sum_i w_i f_i(m)$ over seven factors drives encoding depth,
forgetting, and retrieval, replacing three separately tuned rules. The
weights $w_i$ are fit to a downstream objective (\cref{sec:method:learn}).}
\label{fig:schematic}
\end{figure}

\paragraph{Encoding depth.}
Following levels-of-processing \citep{craik1972lop}, higher-value items
are encoded more deeply. We map the normalised value
$\bar V(m)=V(m)/\sum_i w_i \in[0,1]$ to four encoding tiers
(shallow / semantic / schematic / meta) by thresholds
$\{0.20,0.45,0.70\}$, so deeper encoding is allocated to higher-value
memories at write time.

\paragraph{Forgetting.}
The forget score combines universal time/access dynamics with a
value-driven resistance:
\begin{equation}
\mathrm{forget}(m) \;=\;
\underbrace{(\Delta t)^{0.7}}_{\text{recency}}
\cdot
\underbrace{\tfrac{1}{1+r}}_{\text{usage}}
\cdot
\underbrace{\tfrac{1}{1+\beta\,V(m)}}_{\text{value resistance}},
\label{eq:forget}
\end{equation}
where $\Delta t$ is age in days and $r$ is retrieval count. High value
$V(m)$ multiplicatively resists forgetting, generalising the
emotion-only resistance of prior cognitive memory models to the full
factor set. Items with the highest forget score are dropped first when
the store exceeds budget.

\paragraph{Retrieval.}
At query time the same value contributes to the retrieval rank,
alongside query-specific similarity and mood congruence, so that durable
high-value memories are surfaced preferentially. Because our central
evaluation isolates the budgeted keep/drop decision
(\cref{sec:exp:blind}), we rank candidates directly by $V(m)$ --- using
the regime-specific goal anchor defined there (the question under
\emph{oracle}, the session topic under \emph{blind}) --- and report
retention of the gold set.

\subsection{Learning the weights from task return}
\label{sec:method:learn}

The weights $\wvec$ are \emph{learned}, not hand-set. Let a memory policy
parameterised by $\wvec$ run the
encode\,$\rightarrow$\,\allowbreak forget\,$\rightarrow$\,\allowbreak
retrieve\,$\rightarrow$\,\allowbreak answer
pipeline on a set of training episodes and return a scalar
$\Rtask(\wvec)$ (here, mean gold-evidence retention under a fixed keep
budget; in a fully API-gated setting, downstream QA accuracy). Because
the pipeline contains discrete keep/drop decisions and an external
answerer, $\Rtask$ is non-differentiable in $\wvec$. We therefore
maximise it with a gradient-free stochastic hill-climb
(\cref{alg:learn}), a lightweight stand-in for CMA-ES
\citep{hansen2016cmaes} appropriate to a seven-dimensional search.

\begin{algorithm}[ht!]
\small
\caption{Gradient-free weight learning (A2).}
\label{alg:learn}
\begin{algorithmic}[1]
\State $\wvec \gets \mathbf{1}$ \Comment{uniform init}
\State $R^\star \gets \Rtask(\wvec)$; \; $\sigma \gets \sigma_0$
\For{$t = 1 \dots T$}
  \State $\wvec' \gets \wvec$
  \For{each factor $i$}
    \If{$\mathrm{rand}() < 0.5$}
      $\wvec'_i \gets \max(0,\; \wvec_i + \mathcal{N}(0,\sigma))$
    \EndIf
  \EndFor
  \State $R' \gets \Rtask(\wvec')$
  \If{$R' > R^\star$} $\wvec \gets \wvec'$; \; $R^\star \gets R'$ \EndIf
  \State $\sigma \gets \gamma\,\sigma$ \Comment{anneal}
\EndFor
\State \Return $\wvec$
\end{algorithmic}
\end{algorithm}

The objective is evaluated on a held-in training split and the learned
$\wvec$ is reported on a disjoint test split; best-so-far return is
monotone non-decreasing by construction. To keep the learned weights
honest and interpretable, we restrict the search to the factors actually
populated by the annotator in a given setting --- otherwise the
optimiser assigns arbitrary non-zero weight to factors that are
identically $0$ and therefore multiply out of \cref{eq:value}.
